Eine Firma will ihr kränkelndes Produkt mit einer
Social-Media-Marketing-Initative retten.
Dazu sollen einige Influencer gesponsort werden, die dann das Produkt
ihren Followern bekannt machen sollen.
Da nicht mehr allzu viel Geld vorhanden ist, ist die Anzahl der
Influencer, die unter Vertrag genommen werden können, limitiert.
Es müssen daher möglichst schnell diese Influencer
auf eine Art bestimmt werden, dass jede Person der Zielgruppe erreicht wird.
Die mit der Auswahl betraute Beratungsfirma mahnt, dass dieses
Problem nicht so leicht zu lösen sei.
Warum?

\begin{loesung}
Das Problem \textit{MARKETING} ist äquivalent zum \textit{SET-COVERING}.
Um dies einzusehen, konstruieren wir eine Eins-zu-eins-Reduktioin wie
folgt.
Sei $J=\{1,\dots,n\}$ die Menge aller Influencer,
die zusammen die Zielgruppe $Z$ erreichen können.
Sei weiter $S_j$, $1\le j\le n$ die Menge der Personen, die Influencer $j$
erreichen kann.
Gesucht wird eine Auswahl $S_{i_1},\dots,S_{i_k}$ von $k$ Influencern
derart, dass die ausgewählten Influencer ebenfalls alle Personen der
Zielgruppe
\[
\bigcup_{j=1}^n S_j
= 
\bigcup_{i=1}^k S_{j_i}
\]
erreichen können.
\begin{center}
\begin{tabular}{r>{$}c<{$}l}
\textit{MARKETING}&\leftrightarrow&\textit{SET-COVERING} \\
\hline
Influencer&\leftrightarrow&Nummern $1\le j\le n$ der Mengen \\
Menge der vom Influencer erreichten Personen&\leftrightarrow&Menge $S_j$ \\
Auswahl der Influencer&\leftrightarrow&Unterfamilie $(S_{j_i})_{1\le i\le k}$ \\
Bedingung ganze Zielgruppe erreicht&\leftrightarrow&
$\bigcup_{j=1}^n S_j = \bigcup_{i=1}^k S_{j_i}$
\end{tabular}
\end{center}
Da \textit{SET-COVERING} ein NP-vollständiges Problem ist, ist dafür kein 
polynomieller Algorithmus bekannt und damit das Problem tatsächlich 
für grosse Zielgruppen nur sehr schwer lösbar.
\end{loesung}

\begin{bewertung}
Reduktionsprinzip ({\bf R}) 1 Punkt,
Wahl eines geeigneten Vergleichsprobleme wählen ({\bf V}) 1 Punkt,
Abbildung für Influencer ({\bf I}) 1 Punkt,
Menge der erreichten Personen ({\bf S}) 1 Punkt,
und Zielgruppenbedingung ({\bf Z}) 1 Punkt.
NP-Vollständigkeit und Konsequenzen ({\bf N}) 1 Punkt.
\end{bewertung}

